\chapter{核心功能实现}

\section{基于角色的访问控制与身份认证实现}

本系统面向多角色协同场景，若缺乏清晰的认证授权机制，容易出现数据越权访问和职责冲突问题。因此，系统在安全设计上采用“登录认证 + 令牌校验 + 角色授权 + 业务再校验”的分层控制策略。用户通过账号密码完成认证后，后端生成 JWT，令牌中包含用户标识和角色声明；后续请求到达后，由安全过滤器解析令牌并恢复身份上下文。接口级权限控制根据角色限制访问入口，服务层再结合业务状态进行二次校验。

该机制的工程价值主要体现在两方面。其一，采用 JWT 的无状态认证方式，可以降低服务端会话维护成本，更适合前后端分离架构；其二，基于角色的访问控制使系统能够在统一平台上兼容多类用户，而无需为每类角色单独构建隔离系统。以订单处理为例，顾客只能创建和查询订单，销售管理员只能审核订单，仓库管理员只能在审核通过后处理出入库，角色边界清晰且具备可验证性。

\section{订单—仓储—生产—质检协同流程实现}

订单—仓储—生产—质检协同流程是本系统最具代表性的闭环业务链路。系统并未将订单处理、库存管理和生产执行视为彼此独立的模块，而是通过订单状态、库存状态和任务状态的联动，使其形成连续的流程驱动机制。顾客下单后，订单进入销售审核；审核通过后，仓库根据库存情况决定是直接发货还是触发生产；生产完成后进入质检；质检合格后再由仓库入库并完成发货。

该流程的实现难点在于状态控制与跨模块一致性。为解决该问题，系统采用服务层集中控制状态流转的方法，将订单、库存、生产和质检相关更新置于统一事务边界内处理。这样既避免了状态跳转混乱，也确保了例如“未质检产品不可入库”“库存不足才可触发生产”等规则能够被严格执行。该设计体现了以流程驱动替代人工口头协同的工程思路。

\section{采购—供应商—仓库协同流程实现}

采购协同流程的目标是保证原材料供应连续性，并使企业内部采购部门与外部供应商之间的信息传递更加及时。系统在库存预警、周计划计算或人工申请触发后生成采购需求，采购管理员确认后形成采购订单。随后，系统向供应商推送待处理通知，供应商登录后执行接单与发货反馈。货物到达企业后，由仓库管理员完成收货入库，并视需要结合质检管理员完成来料检验。

该流程的实现重点在于多主体状态同步。采购订单在创建、确认、发货、到货和入库等节点均会更新状态，并通过消息通知同步给相关角色。这样可显著降低电话、即时通信软件等非结构化沟通方式造成的信息遗漏风险。同时，系统保留各处理节点记录，为后续供应商履约评估和采购追踪提供数据基础。

\section{库存预警与周计划自动生成算法实现}

库存预警与周计划生成的设计目标并非追求复杂的预测模型，而是在企业现有业务数据基础上构建可落地的规则计算机制。库存预警首先针对成品和原材料分别计算预计可用库存，当预计库存低于安全库存阈值时即触发预警。对于产品 $p$ 在某一计划周期内的预计库存，可表示为

\begin{equation}
I_{p}^{\mathrm{proj}} = I_{p}^{\mathrm{cur}} + I_{p}^{\mathrm{in}} - D_{p}^{\mathrm{ord}} - D_{p}^{\mathrm{plan}}
\end{equation}

其中，$I_{p}^{\mathrm{cur}}$ 表示当前库存，$I_{p}^{\mathrm{in}}$ 表示在途入库数量，$D_{p}^{\mathrm{ord}}$ 表示已确认订单需求，$D_{p}^{\mathrm{plan}}$ 表示计划期间额外消耗量。当 $I_{p}^{\mathrm{proj}}$ 小于安全库存时，系统生成预警信息。

在周计划生成阶段，系统进一步综合需求量、安全库存和现有供给，计算建议生产量。其基本思想可表示为

\begin{equation}
Q_{p}^{(w)} = \max \left(0, \hat{D}_{p}^{(w)} + S_{p}^{(w)} - I_{p}^{\mathrm{cur}} - I_{p}^{\mathrm{in}} \right)
\end{equation}

其中，$\hat{D}_{p}^{(w)}$ 表示未来一周需求预测值，$S_{p}^{(w)}$ 表示安全库存目标，$Q_{p}^{(w)}$ 表示建议生产量。若产品存在物料清单，则可进一步估算原材料需求量

\begin{equation}
M_{m}^{(w)} = \sum_{p=1}^{n} Q_{p}^{(w)} \times B_{pm}
\end{equation}

其中，$B_{pm}$ 表示产品 $p$ 对原材料 $m$ 的单位消耗量。通过上述规则，系统能够生成生产周计划与采购周计划建议。其优点在于实现成本低、可解释性强、便于业务人员理解和修正。

为了保证规则计算流程清晰，系统在计划服务中采用如下处理步骤。

\begin{algorithm}[H]
    \KwIn{订单需求集合、库存数据、在途数据、安全库存阈值、BOM 数据}
    \KwOut{生产周计划、采购周计划}
    汇总未来一周已确认订单需求\;
    计算产品预计可用库存\;
    \ForEach{产品}{
        若预计库存低于安全库存，则计算建议生产量\;
    }
    根据建议生产量和 BOM 关系换算原材料需求\;
    \ForEach{原材料}{
        若预计物料库存低于阈值，则生成采购建议量\;
    }
    写入生产周计划与采购周计划记录\;
    向生产管理员和采购管理员发送计划通知\;
    \caption{周计划自动生成流程}
\end{algorithm}

\section{实时消息推送与状态同步实现}

在多角色协同系统中，若业务状态变化不能及时传递给相关岗位，即使系统具备完整功能，也难以真正提升协同效率。为此，本系统采用 WebSocket/STOMP 构建实时消息机制，将订单审核、库存预警、采购下发、生产任务生成和质检结果等关键事件转换为可订阅消息。后端以业务事件为触发点向指定用户或角色主题推送通知，前端在收到消息后更新页面待办项和状态展示。

该机制的关键价值在于将原本被动查询式的信息获取转变为主动推送式的状态同步。例如，销售审核通过后，仓库无需周期性刷新页面等待新任务；质检完成后，仓库可立即获知是否允许入库；采购订单生成后，供应商可以直接接收待处理提醒。由此可见，实时消息机制不仅改善了用户体验，更增强了业务流程的连续性。

\section{质量追溯功能实现}

质量追溯功能的核心在于将订单、生产、质检和库存流转记录进行贯通，并围绕统一的业务标识建立查询链路。系统在订单生成时即建立基础追踪主线，在后续生产任务、质检结果和出入库记录中保留与订单或产品批次相关的关联字段。当顾客或管理人员发起追溯查询时，后端服务依据订单号或业务标识聚合相关记录，并按时间顺序返回处理链路。

从工程实现角度看，质量追溯并不依赖额外复杂模型，而是依赖业务处理过程中对关键数据节点的完整留痕。只要系统能够保证生产记录、质检记录和库存记录的时序性与关联性，就能够在问题发生后快速定位处理节点和责任环节。因此，质量追溯既是系统透明度的重要体现，也是后续质量管理优化的数据基础。